\subsection{Marco Teórico}

Este estudio reúne conceptos de la teoría del aprendizaje basado en competencias, la analítica del aprendizaje y la minería de datos educativos para comprender cómo la información académica puede aprovecharse en el fortalecimiento de las competencias y en la toma de decisiones dentro de la educación superior. Desde la perspectiva del aprendizaje por competencias, el desempeño académico no se limita a la adquisición de conocimientos, sino que también considera la capacidad del estudiante para aplicar lo aprendido en situaciones reales. Según lo expuesto por \cite{tobon2013formacion} en 2013, las competencias integran el saber, el saber hacer y el saber ser, constituyendo uno de los pilares de los procesos actuales de evaluación educativa. Por otra parte, la analítica del aprendizaje permite a las instituciones identificar patrones asociados al rendimiento, al riesgo académico y a las oportunidades de mejora a partir del análisis de grandes volúmenes de información educativa. En este contexto, las técnicas de minería de datos educativos facilitan la construcción de modelos predictivos y herramientas de diagnóstico que convierten los datos académicos en información útil para la toma de decisiones. De acuerdo con \cite{cobo2011aprendizaje}, el uso de herramientas analíticas contribuye a respaldar las decisiones institucionales con evidencia objetiva y favorece el mejoramiento continuo de la calidad educativa mediante una gestión más efectiva del conocimiento. 
Desde una perspectiva teórica más amplia, la teoría sociocognitiva de Bandura (1977) y el constructivismo de Vygotsky (1978) permiten comprender el aprendizaje como un proceso que surge de la interacción entre el estudiante, su entorno y los recursos educativos disponibles \cite{bandura1977self} \cite{vygotsky1978mind}. Ambos enfoques coinciden en que el aprendizaje se construye de manera activa y está influenciado por factores sociales y contextuales. Esta visión respalda el enfoque del sistema de información propuesto, en el que los datos no solo se recopilan y almacenan, sino que también se analizan considerando las condiciones individuales y colectivas que caracterizan a cada estudiante.  En conjunto, estos fundamentos teóricos permiten abordar el aprendizaje desde diferentes perspectivas complementarias: la medición de competencias, el análisis de datos educativos y la comprensión del contexto en el que ocurre el proceso formativo. Bajo esta visión, el sistema de información propuesto busca no solo estimar el desempeño entre las pruebas Saber 11 y Saber Pro, sino también generar información útil para apoyar la toma de decisiones institucionales, fortalecer el desarrollo de competencias y contribuir a la reducción de brechas académicas.


\subsection{Marco Conceptual}

Explorando el contexto de la investigación en Colombia, la evaluación de las competencias que presenta un
estudiante se estructura en torno a los instrumentos estandarizados diseñados por el Instituto Colombiano
para la Evaluación de la Educación (ICFES). La prueba Saber~11, aplicada al finalizar la educación media
o bachillerato, permite medir las habilidades y conocimientos que los estudiantes han adquirido en áreas
como Lectura crítica, Matemáticas, Ciencias naturales, Sociales y ciudadanas, e Inglés \cite{icfes2024_informeSaber11_2023}. Por su
parte, la prueba Saber~Pro evalúa el desarrollo de competencias genéricas y específicas alcanzadas durante
la formación universitaria, proporcionando un marco de referencia que funciona a nivel nacional para valorar
la calidad del proceso formativo de las instituciones de educación superior \cite{sapiencia2022_infographicBestSaberPro}. Ambas pruebas, al ser
de carácter obligatorio y comparables en el tiempo, constituyen una fuente de información esencial para
estudiar la continuidad y/o evolución de las habilidades académicas a lo largo del proceso educativo.

Existe un fenómeno conocido como ``\textbf{valor agregado educativo}'', concepto introducido en la literatura de eficacia escolar por \cite{hanushek1971teacher} (1971) y adoptado en Colombia por el Ministerio de Educación Nacional y el ICFES como indicador del progreso o ganancia en el nivel de competencias que experimenta el estudiante entre dos momentos de su trayectoria formativa, y es precisamente el análisis de los resultados obtenidos en estas pruebas el que permite comprender este valor agregado \cite{icfes2021valoragregado}. La investigación realizada en Colombia por \cite{maestre2022_valorAgregadoCalidadEducativa}, en la Universidad del Norte, plantea que este indicador permite
estimar el aporte que realiza la institución al aprendizaje del estudiante, considerando su punto de partida
y su contexto socioeconómico. Dicho valor no mide únicamente resultados finales, sino la diferencia entre
el desempeño alcanzado y aquel que sería esperado según las condiciones iniciales.

La medición del valor agregado se apoya en modelos predictivos que estiman el puntaje esperado en
Saber~Pro a partir de variables como los resultados de Saber~11, los antecedentes académicos y los factores
contextuales del estudiante. Por ejemplo, el estudio ``Predicción rendimiento estudiantes pruebas Saber Pro
en pandemia junto con las características socioeconómicas'' propone el uso de algoritmos de minería de
datos ---como árboles de decisión y regresión logística--- para generar anticipaciones al desempeño final;
de esta forma, la diferencia entre el puntaje real y el valor pronosticado constituye una medida cuantitativa
del valor agregado \cite{vanegas2022_prediccionSaberPro}.

Es importante delimitar los factores a examinar; así, el modelo conceptual considera la relación entre tres
grandes conjuntos de variables: las competencias iniciales, representadas por los puntajes en Saber~11;
las competencias terminales, representadas por los puntajes en Saber~Pro; y los factores socioeconómicos
y de contexto que influyen en dicha relación. Este estudio considera que las competencias evaluadas en la prueba Saber 11 representan el punto de partida sobre el cual los estudiantes continúan desarrollando sus habilidades durante la educación superior. En consecuencia, quienes obtienen mejores resultados en esta prueba suelen contar con una base académica que favorece su desempeño posterior en Saber Pro. Sin embargo, esta relación no depende únicamente de las capacidades iniciales del estudiante, sino también de diversos factores sociales, económicos e institucionales que influyen en su trayectoria formativa.
Entre los aspectos contextuales más relevantes se encuentran el nivel socioeconómico, la escolaridad de los padres, el tipo de institución educativa de procedencia, el acceso a recursos tecnológicos y la ubicación geográfica. En general, los estudiantes que provienen de entornos con mayores recursos educativos y económicos suelen disponer de más oportunidades de aprendizaje, lo que puede reflejarse en mejores resultados en las pruebas estandarizadas. De igual forma, quienes han tenido acceso a instituciones educativas con mayores recursos o a herramientas tecnológicas de apoyo suelen mostrar desempeños más favorables en Saber 11 y una evolución académica más consistente durante sus estudios universitarios.

A lo largo de la formación profesional, la universidad desempeña un papel fundamental en el desarrollo de estas competencias. En este proceso intervienen factores asociados a la Institución de Educación Superior (IES), como la calidad del programa académico, la pertinencia de los contenidos curriculares, las estrategias pedagógicas utilizadas, la disponibilidad de recursos tecnológicos y el desempeño académico alcanzado por los estudiantes, medido a través de indicadores como el promedio acumulado (GPA) o el número de asignaturas reprobadas. 

Estos elementos reflejan la capacidad de la institución para fortalecer las competencias adquiridas durante la educación media y contribuir al desarrollo de habilidades profesionales de mayor complejidad. Diversos estudios han encontrado que las instituciones con procesos consolidados de acreditación y gestión de la calidad suelen generar mayores niveles de valor agregado en las competencias evaluadas, lo que indica que la calidad institucional tiene una influencia importante en la evolución académica de los estudiantes \cite{maestre2022_valorAgregadoCalidadEducativa}.

En cuanto a los conceptos metodológicos del trabajo, la regresión lineal es uno de los modelos predictivos
más utilizados en el aprendizaje automático y la estadística, cuyo objetivo es describir la relación entre
una variable dependiente continua y una o más variables independientes mediante una función lineal. En
su forma simple, busca minimizar la suma de los errores cuadráticos entre los valores observados y los
estimados, lo que la convierte en una herramienta eficaz para evaluar tendencias entre factores como las
competencias académicas iniciales y los resultados finales \cite{google2025linear}. Este método resulta particularmente
útil cuando se requiere interpretar el peso o la influencia de cada variable predictora en el resultado,
permitiendo entender cómo ciertos indicadores de Saber~11 pueden anticipar el rendimiento en Saber~Pro
\cite{google2025linear}.

Los árboles de decisión representan un método no paramétrico que divide recursivamente el espacio de
los datos en regiones homogéneas según los valores de las variables predictoras. Cada nodo de decisión
contiene una regla condicional, y las hojas finales representan los valores de salida estimados. Su principal
ventaja radica en la capacidad de manejar relaciones no lineales y de ofrecer una interpretación intuitiva
mediante reglas ``si-entonces'' que facilitan la comprensión del modelo por parte de actores institucionales
\cite{mienye2023survey}. Además, pueden extenderse a técnicas como los bosques aleatorios (\textit{Random Forests})
o los \textit{Gradient Boosted Trees}, que mejoran su rendimiento reduciendo la varianza de los modelos
individuales \cite{mienye2023survey}.

El algoritmo de los $k$ vecinos más cercanos (KNN) se basa en la proximidad o similitud entre
observaciones. Este enfoque asume que las instancias cercanas en el espacio de características tienden a
compartir valores similares de la variable objetivo. Para una nueva muestra, el modelo calcula las distancias
a todas las observaciones existentes y asigna la predicción en función del promedio o la moda de los $k$
vecinos más cercanos \cite{bhatia2010survey}. KNN no requiere un entrenamiento paramétrico complejo, lo que lo
convierte en una técnica flexible y útil en contextos educativos donde el número de variables y patrones
puede ser diverso \cite{bhatia2010survey}.

Las redes neuronales artificiales (\textit{RNA}), y en particular el perceptrón multicapa (\textit{MLP}, por sus siglas en inglés), constituyen modelos de aprendizaje supervisado capaces de aproximar relaciones no lineales complejas entre las variables de entrada y la variable objetivo. Un \textit{MLP} organiza el cálculo en capas sucesivas de neuronas: cada neurona combina linealmente las salidas de la capa anterior y aplica una función de activación no lineal, la función de activación usada más habitualmente es la Unidad Lineal Rectificada (\textit{Rectified Linear Unit} o \textit{ReLU}), lo que permite a la red representar interacciones que un modelo lineal no captura. En problemas de regresión, la capa de salida emplea una activación lineal y el entrenamiento ajusta los pesos minimizando una función de pérdida, como el error cuadrático medio, mediante descenso de gradiente y retropropagación \cite{goodfellow2016deep}. Su principal ventaja en este trabajo radica en la capacidad de modelar conjuntamente los distintos predictores cuando se dispone de un volumen de datos suficientemente grande, condición que se satisface en el conjunto emparejado utilizado.

Por su parte, las técnicas de segmentación o agrupamiento (\textit{clustering}) permiten organizar a los estudiantes en grupos homogéneos según la similitud de su perfil de entrada, sin emplear la variable objetivo durante ese proceso. El algoritmo \textit{K-means} es una de las técnicas de agrupamiento más utilizadas. Su funcionamiento consiste en organizar las observaciones en $K$ grupos, de manera que cada una quede asociada al centroide más cercano dentro de su grupo \cite{mcqueen1967some}. Aunque K-means no realiza predicciones de forma directa, puede utilizarse como base para construir modelos predictivos mediante una estrategia conocida como \textit{cluster-then-predict} o agrupar y predecir. En una primera etapa, los estudiantes se agrupan según las características de sus perfiles de entrada. Posteriormente, para cada grupo se calcula una respuesta asociada a la variable objetivo, que puede corresponder al promedio de los resultados observados o a una regresión lineal ajustada específicamente para ese grupo. Cuando se incorpora un nuevo estudiante, la predicción se obtiene asignándolo al grupo cuyos integrantes presentan características más similares y aplicando la estimación definida para dicho grupo. Este enfoque parte del supuesto de que estudiantes con perfiles iniciales similares tienden a presentar desempeños futuros similares.

Dentro del proceso de evaluación de modelos predictivos, la Raíz del Error Cuadrático Medio (\textit{RMSE}, por sus siglas en inglés) es una métrica ampliamente empleada para medir la precisión del modelo. Se define como la raíz cuadrada de la media de los errores al cuadrado, proporcionando una medida del error promedio en las mismas unidades que la variable dependiente \cite{hodson2022rmse}. Cuanto menor sea el valor del RMSE, mayor será la precisión del modelo. En estudios orientados a la predicción del desempeño académico, esta métrica permite medir de manera directa qué tan cerca se encuentran las estimaciones realizadas por el modelo de los resultados reales obtenidos por los estudiantes \cite{hodson2022rmse}. 

Por su parte, el coeficiente de determinación ($R^{2}$) permite evaluar qué proporción de la variabilidad observada en los datos puede explicarse mediante el modelo. Su valor oscila entre 0 y 1, donde los valores más cercanos a 1 indican una mayor capacidad para representar el comportamiento de la variable analizada. Esta métrica resulta útil para evaluar la calidad tanto de los modelos lineales como de los no lineales, ya que permite identificar en qué medida el rendimiento académico está asociado con los factores incluidos en el análisis \cite{chicco2021r2}. 

El preprocesamiento de los datos constituye una etapa fundamental para garantizar resultados confiables. Durante esta fase se realizan actividades como el tratamiento de valores faltantes, la transformación de variables categóricas y la adecuación de las escalas de medición. Diversos estudios señalan que una preparación adecuada de los datos contribuye a mejorar el desempeño de los modelos y a disminuir la posibilidad de obtener resultados sesgados \cite{han2022data}. 
Asimismo, herramientas como Pandas, NumPy y Scikit-learn facilitan la automatización de estas tareas dentro del entorno de Python, favoreciendo la consistencia del proceso y la reproducibilidad de los análisis \cite{han2022data}.

Dentro de esta etapa, la normalización de las variables requiere especial atención, ya que existen diferentes métodos y cada uno puede influir de manera distinta en el comportamiento de los modelos. La normalización Min-Max transforma los valores de cada variable a un rango determinado, generalmente entre $[0, 1]$, manteniendo la forma original de la distribución. Por otro lado, la estandarización o puntaje z ajusta los datos para que tengan una media igual a cero y una desviación estándar unitaria. Este procedimiento suele ser especialmente útil en modelos cuyo desempeño depende de la escala de las variables de entrada \cite{gil2014analisis}.

En el contexto del aprendizaje automático aplicado a la educación, el campo de \textit{Educational Data Mining} (EDM) ha demostrado que los algoritmos supervisados como la regresión, los árboles de decisión y KNN son útiles para predecir el rendimiento de los estudiantes y detectar factores de riesgo académico
\cite{ozyurt2023educational}. Diversos estudios han resaltado la importancia de combinar estos modelos con métricas de
evaluación objetivas y sistemas de visualización interactivos para facilitar la interpretación de los resultados
por parte de los tomadores de decisiones institucionales \cite{ozyurt2023educational}.

A partir de estos conceptos y criterios en conjunto, la visualización y los sistemas de información educativos permiten traducir los resultados de los modelos predictivos en herramientas accesibles para la gestión académica. Los tableros de control (\textit{dashboards}) y reportes automatizados facilitan la comprensión de patrones de rendimiento, fomentando la toma de decisiones basada en datos dado que reflejan dinámicas y tendencias ocultas, que brindan mayor información y precisión al momento de tomar decisiones \cite{alamri2021explainable}. La implementación de un sistema de información predictivo y comparativo permitirá optimizar el análisis del desempeño académico institucional, fortaleciendo la planeación y las estrategias de mejora \cite{alamri2021explainable}.

De este modo, se observa que la educación es un proceso acumulativo en el que las competencias previas,
las condiciones contextuales y las oportunidades institucionales interactúan para definir el desarrollo
académico del estudiante. Por ello, la comprensión de estas interrelaciones es esencial para diseñar
estrategias de intervención educativa, orientar políticas institucionales y fortalecer los procesos de
formación, evaluación y seguimiento de los resultados de aprendizaje.

\subsection{Estado del Arte}

En los últimos años, diferentes investigaciones realizadas en Colombia han estudiado la relación entre los resultados obtenidos en las pruebas Saber 11 y Saber Pro, así como la posibilidad de predecir el desempeño de los estudiantes en esta última evaluación. En 2024, el estudio de \cite{mendieta2024modelo} desarrolló un modelo descriptivo y predictivo orientado a apoyar el mejoramiento de los resultados en Saber Pro. Para ello, utilizó información histórica y variables de caracterización estudiantil con el propósito de identificar los factores que ejercen mayor influencia sobre el rendimiento académico. De manera similar, \cite{cervera2021modelo} presentó en 2021 un análisis basado en algoritmos como Naive Bayes y \textit{k-Nearest Neighbors} para predecir los resultados de Saber Pro. Los hallazgos mostraron que variables relacionadas con el contexto familiar y social, como el nivel educativo de la madre, la composición del hogar, el lugar de residencia y la disponibilidad de recursos en el entorno doméstico, tienen una influencia importante en el desempeño de los estudiantes.

Por otra parte, \cite{granados2025modelado} aplicó en 2025 modelos de regresión logística multinomial para estimar los diferentes niveles de desempeño en Saber Pro a partir de variables académicas y contextuales. Asimismo, \cite{del2022indice} propuso en 2022 metodologías de seguimiento longitudinal entre las pruebas Saber 11 y Saber Pro con el objetivo de medir el valor agregado generado durante la formación universitaria. Este trabajo destacó la relevancia de factores como las condiciones institucionales, el contexto geográfico y los mecanismos de admisión en la explicación de los resultados obtenidos. En conjunto, estas investigaciones muestran que existe evidencia suficiente sobre el potencial de los datos de Saber 11 para anticipar el desempeño futuro de los estudiantes cuando se complementan con variables académicas y de contexto. Estos antecedentes respaldan la viabilidad del sistema de información propuesto y refuerzan la importancia de desarrollar herramientas analíticas que permitan comprender mejor la evolución de las competencias y apoyar los procesos de mejora institucionales.

